\section{Optional GPS-Derived Trip-Attribute Observations}
\label{sec:gps-trip-attribute-observations}

Boarding and alighting counts are the primary observations in this model.  A
case study may also have a sample of complete passenger journeys reconstructed
from GPS data.  Such a sample provides information about the distribution of
journey attributes, but it is not, by itself, an additional OD table or a
calibrated demand count.  It can therefore be included as an optional,
auxiliary observation model without changing the count-assignment model.

\subsection{Data description}

One row in the auxiliary data represents one reconstructed public-transport
journey, from its first boarding to its final alighting.  A representative
schema contains the following fields:

\begin{description}
  \item[\texttt{departure\_hour}] the local clock hour of the first boarding,
    obtained by taking the floor of the first-boarding timestamp;
  \item[\texttt{travel\_time}] the sum of the in-vehicle travel times of the
    journey's legs, measured in seconds;
  \item[\texttt{airline\_distance}] the origin--destination air-line distance,
    measured in metres.  The geometric points and distance formula must be
    declared by the case study;
  \item[\texttt{traveled\_stops}] the number of stops traversed by the
    journey;
  \item[\texttt{transfers}] the number of changes of vehicle.
\end{description}

For the current case study, the GPS sample covers 52 weekdays between
3~October and 23~December~2022 under reported normal-traffic conditions.  It
has no sampling or expansion weights, and it is not assumed to represent the
absolute number of journeys in each departure hour.  Its intended information
content is the shape of the trip-time and transfer distributions.  Multiple
rows can come from the same individual.  Moreover, the APC observations are
exhaustive, so GPS journeys recorded on 22~November~2022 also occur in the APC
data for that day.  The two sources must therefore not be treated as fully
independent replicate samples.

The absence of GPS records in the 02:00--03:00 hours is a service fact rather
than missing data: no vehicle operates in that interval.  The 03:00--04:00
interval has only very sparse service and should not be assigned its own
distribution without a documented pooling rule.  The row count is version
dependent and is recorded by the data audit together with the file checksum,
rather than being a fixed mathematical constant in this report.

The current journey reconstruction removes trips whose initial, transfer, or
final stops cannot be deduced; merges consecutive legs using the same line and
direction when the intervening transfer time is below 120 seconds; and slices
records that incorrectly merge a same-line opposite-direction return or a
transfer longer than 20 minutes.  An overnight journey with an implausibly
long duration and zero-distance records with zero traversed stops are treated
as invalid.  Other long journeys and zero-distance journeys with non-zero
stops require a documented case-owner decision because they may be distinct
journeys merged by the reconstruction algorithm.  Every such decision belongs
in the data-cleaning audit, not in the likelihood definition.

In particular, the trip time in this schema is not door-to-door journey time:
it excludes access walking, egress walking, initial waiting, and transfer
waiting unless an explicitly extended data contract supplies those quantities.
The routing representation used by the model must compute the same quantity.
Adding waiting time to the GPS field without adding the corresponding quantity
to the route attributes would create a systematic measurement mismatch.

The file may also contain exact timestamps, origin and destination identifiers,
line or trip identifiers, or visited-stop sequences.  These richer fields are
useful for validation and can support a more detailed observation model, but
the aggregate formulation below does not require them.  If those identifiers
are unavailable, the GPS data constrain aggregate journey attributes only; they
cannot identify individual OD cells.

\subsection{Data contract and assumptions}

The optional data term is valid only after a case-specific data audit.  The
following assumptions must be stated and checked:

\begin{enumerate}
  \item The units, time zone, journey boundaries, stop definition, transfer
    definition, and distance definition agree with the corresponding quantities
    produced by the fixed-routing model.
  \item The sample is interpreted as information about the distribution of
    journey attributes within a declared departure-time stratum, not as an
    estimate of the absolute passenger volume in that stratum.  Sampling weights
    are used if they exist; otherwise the likelihood must not treat the raw row
    count as a calibrated demand total.
  \item Repeated journeys by the same person, and overlap with an exhaustive
    APC data set, are recorded.  The GPS rows are therefore not automatically
    conditionally independent of one another or of the count observations.
    Clustering, an effective sample size, or an explicit likelihood tempering
    factor is required when the term is used for estimation.
  \item Departure-time strata are chosen with enough observations to support a
    stable distribution.  A separate distribution for every clock hour is not
    required.  Hours with no service contribute no GPS likelihood term, and a
    sparse hour may be pooled with neighbouring hours under a documented rule.
  \item Invalid or ambiguous journeys are removed only under a documented
    data-owner rule.  The original file, the cleaning rule, the removed-row
    identifiers, and the resulting checksum are retained.  An observed GPS bin
    with no modelled route support is a path-attribute or data-contract failure,
    not an observed zero that may be discarded silently.
\end{enumerate}

Because GPS coverage can vary by person, day, line, or time of day, the
assumption that the sample is representative must be treated as a modelling
choice and tested.  A distributional term should be omitted, or used only for
diagnostics, when that assumption cannot be defended.

\subsection{Attribute-response construction}

Let $j$ index the free OD--time cells and let $x_j(\theta)$ be the demand
generated by the selected demand model.  The existing count response is
\begin{equation}
  \boldsymbol{\mu}_{\mathrm{count}}(\theta)
  =\rho\bigl(A\boldsymbol{x}(\theta)+\boldsymbol{c}\bigr),
\end{equation}
where $A$ is the fixed-routing measurement operator and $\boldsymbol{c}$ is
the contribution of fixed positive demand.

For the GPS data, let $s$ index a declared departure-time stratum and let
$b$ index a bin of one attribute or of a declared joint attribute vector (for
example, in-vehicle travel-time bins crossed with transfer counts).  Define
$B_{sb,j}$ as the route-share-weighted probability that demand cell $j$
produces an assigned journey in stratum $s$ and bin $b$.  The predicted GPS
mass is
\begin{equation}
  \nu_{sb}(\theta)
  =\sum_j B_{sb,j}x_j(\theta)+\nu^{\mathrm{fixed}}_{sb}.
  \label{eq:gps-attribute-mass}
\end{equation}
The fixed-demand term is included for the same reason as the fixed count
offset.  Conditional on stratum $s$, the predicted attribute distribution is
\begin{equation}
  p_{sb}(\theta)
  =\frac{\nu_{sb}(\theta)}{\sum_{b'}\nu_{sb'}(\theta)}.
  \label{eq:gps-attribute-probability}
\end{equation}
The operator $B$ is a second view of the same fixed routes; it does not create
new route choices or alter $A$.

In the public implementation, the route-level contributions are represented
by a \texttt{GravityRouteShare} for each free OD--time cell and each assigned
route.  The corresponding
\texttt{GravityAttributeResponseOperator} stores the resulting category
contributions either as a small dense reference matrix or as contiguous
category shards.  Its fixed-demand offset is kept separately, and its
forward/adjoint products have the same compact free-OD dimension as the count
operator.  The attribute-response artifact has its own provenance fingerprint
so that changing attribute bins invalidates $B$ without rebuilding $A$.
Support is checked before the likelihood is evaluated: a positive observed
count in a category with no free-route or fixed-demand support is a strict
preflight failure, whereas a zero count in such a category is not silently
interpreted as a route observation.

\subsection{Optional GPS likelihood and MAP objective}

Let $n_{sb}$ be the observed GPS histogram and
$N_s=\sum_b n_{sb}$.  If the sample is treated as a representative sample
within each stratum, the basic auxiliary likelihood is
\begin{equation}
  \boldsymbol{n}_s\mid\theta
  \sim \operatorname{Multinomial}
       \bigl(N_s,\boldsymbol{p}_s(\theta)\bigr).
  \label{eq:gps-multinomial}
\end{equation}

The multinomial form is selected because the GPS file supplies a sample of
completed journeys but no expansion weights or known capture rate.  Conditional
on $N_s$, it uses the information that is actually identified---the relative
frequency of journey attributes within a stratum---without interpreting the
number of GPS records as an absolute passenger volume.  A Poisson likelihood on
$n_{sb}$ would require a calibrated exposure or sampling rate; without one it
would incorrectly add a second, arbitrary production measurement to the
exhaustive APC counts.

The ordinary multinomial nevertheless treats all rows as independent.  That is
not credible here because several rows may come from one individual, the sample
covers multiple days, journey reconstruction is uncertain, and some GPS
journeys overlap the APC observations.  When a cluster or uncertainty estimate
is available, a Dirichlet--multinomial model is preferred:
\begin{equation}
  \boldsymbol{n}_s\mid\theta
  \sim \operatorname{DirichletMultinomial}
       \bigl(N_s,\kappa_s\boldsymbol{p}_s(\theta)\bigr),
  \label{eq:gps-dirichlet-multinomial}
\end{equation}
where the concentration $\kappa_s$ controls the extra-multinomial variation.
Large concentration approaches the ordinary multinomial; smaller concentration
represents the loss of information caused by clustering and reconstruction or
day-to-day variation.  An equivalent practical route is to use a
cluster-adjusted effective sample size, but the adjustment must be documented.

If neither a cluster estimate nor an effective sample size is available, a
tempered multinomial log-likelihood is a conservative composite-likelihood
approximation:
\begin{equation}
  \ell_{\mathrm{GPS,weighted}}(\theta)
  =\lambda_{\mathrm{GPS}}\ell_{\mathrm{GPS}}(\theta),
  \qquad 0\leq\lambda_{\mathrm{GPS}}\leq 1.
  \label{eq:gps-tempered-likelihood}
\end{equation}
The tempering factor prevents the unweighted GPS sample from overwhelming the
primary APC evidence and acknowledges the cross-source dependence.  It is not
mathematically equivalent to a Dirichlet--multinomial distribution and must be
reported as a modelling choice.  A Dirichlet--multinomial concentration or an
effective sample size and a tempering factor should not both be used to reduce
the same evidence unless their separate roles are explicitly justified.

The implementation in
\path{public_transportation.inference.gravity.likelihoods} applies these
models to the route-response masses after normalization separately within
each declared stratum.  It treats a supplied effective sample size as the
Dirichlet concentration in the current aggregate contract, and reports this
choice as part of the uncertainty metadata.  The multinomial,
Dirichlet--multinomial, and tempered forms are mutually exclusive
declarations; the tempering factor is not silently applied on top of a
Dirichlet--multinomial concentration.  All three functions are
JAX-differentiable and the zero-tempering limit contributes exactly zero.  The
count-only gravity objective remains unchanged when no auxiliary channel is
enabled.

With the existing count log-likelihood and prior penalty, the optional MAP
objective is
\begin{equation}
  \mathcal{J}_{\mathrm{MAP}}(\theta)
  =-\ell_{\mathrm{count}}(\theta)
   -\lambda_{\mathrm{GPS}}\ell_{\mathrm{GPS}}(\theta)
   +\mathcal{R}(\theta).
  \label{eq:gps-augmented-map}
\end{equation}
Setting $\lambda_{\mathrm{GPS}}=0$ recovers the existing count-only model.
The GPS term is therefore an optional source of information, not a replacement
for the boarding/alighting likelihood.  In ML estimation $\mathcal{R}$ is
omitted; in MAP estimation it contains the declared prior or regularization.

The implementation connects these terms through
\path{public_transportation.inference.gravity.aggregate_channel}.  One
validated histogram becomes one optional observation channel.  The gravity
objective retains the APC contribution and each auxiliary contribution
separately, then sums them before adding the regularization term.  Both the
batched-forward and adjoint gradients propagate the auxiliary derivative
through the same free-OD demand vector.  With an empty channel bundle, the
objective, gradients, model fingerprint, and checkpoints retain the existing
count-only behavior.

The aggregate GPS sample can constrain shared gravity parameters that change
the mix of short and long journeys or the number of transfers.  Without
origin--destination identifiers, however, it cannot supply independent
information for every OD cell.  Its influence must be reported separately from
the count likelihood and assessed through sensitivity to the departure-time
strata, attribute bins, effective sample size, and
$\lambda_{\mathrm{GPS}}$.

If the auxiliary file is unavailable, fails the data contract, or is judged
non-representative, equations~\eqref{eq:gps-attribute-mass}--\eqref{eq:gps-augmented-map}
are simply omitted and the existing count-based estimation procedure is
unchanged.
